\subsection{约化根系}

称根系是约化的，如果系统的每个根都是不可分根（no. 3）。

命题 12. 假设 R 不可约且约化。

\begin{enumerate}
\def\labelenumi{(\roman{enumi})}
\item
  比值 \(\frac{(\beta|\beta)}{(\alpha|\alpha)}\) 对 \(\alpha \in \mathbb{R}, \beta \in \mathbb{R}\) 必须取 \(1, 2, \frac{1}{2}, 3, \frac{1}{3}\) 中的一个值。
\item
  集合 \((\alpha |\alpha)\)（其中 \(\alpha \in \mathrm{R}\)）至多含有两个元素。
\end{enumerate}

由于 R 不可约，根在 W(R) 下的变换生成 V（no. 2, Cor. of Prop. 5）。因此，对任意根 \(\alpha, \beta\)，存在一个根 \(\beta'\)，使得 \((\alpha|\beta') \neq 0\) 且 \((\beta'|\beta') = (\beta|\beta)\)。由第 1 号的公式 (9) 及第 3 号的列表，\(\frac{(\beta'|\beta')}{(\alpha|\alpha)}\) 取 \(1, 2, \frac{1}{2}, 3, \frac{1}{3}\) 中的一个值（回顾一下，系统假设为约化的）。将 \((x|y)\) 乘以适当的标量，可以假设对于某些根 \((\alpha|\alpha) = 1\)，而 \((\beta|\beta)\) 对 \(\beta \in \mathbb{R}\) 的其他可能值为 2 和 3。2 和 3 不可能同时取得，因为在这种情况下会存在 \(\beta \in \mathbb{R}, \gamma \in \mathbb{R}\)，使得 \(\frac{(\gamma|\gamma)}{(\beta|\beta)} = \frac{3}{2}\)，这与我们上面看到的相矛盾。

命题 13. 假设 R 不可约、非约化且秩 \(\geqslant 2\)。

\begin{enumerate}
\def\labelenumi{(\roman{enumi})}
\item
  不可分根的集合 \(\mathbf{R}_0\) 是 V 中的根系；该系统不可约且约化；并且 \(\mathrm{W}(\mathrm{R}_0) = \mathrm{W}(\mathrm{R})\)。
\item
  设 \(\mathbf{A}\) 是根 \(\alpha\) 的集合，其中 \((\alpha |\alpha)\) 取最小值 \(\lambda\)。则 \(\mathbf{A}\) 中任意两个不成比例的元素都正交。
\item
  设 B 是 \(\beta \in \mathbb{R}\) 的集合，满足 \((\beta |\beta) = 2\lambda\)。则 \(B \neq \emptyset\)，\(R_0 = A \cup B\)，\(R = A \cup B \cup 2A\)。
\end{enumerate}

若 \(\alpha \in \mathbb{R} - \mathbb{R}_0\)，则 \(\frac{1}{2}\alpha \in \mathbb{R}\)，但 \(\frac{1}{2}\left(\frac{1}{2}\alpha\right) \notin \mathbb{R}\)（命题 8），所以 \(\frac{1}{2}\alpha \in \mathbb{R}_0\)。这证明了 \(\mathbb{R}_0\) 满足 \((\mathrm{RS}_1)\)。显然，对所有 \(\alpha \in \mathbb{R}\)，\(s_{\alpha, \alpha'}(\mathbb{R}_0) = \mathbb{R}_0\)，所以 \(\mathbb{R}_0\) 满足 \((\mathrm{RS}_{II})\) 和 \((\mathrm{RS}_{III})\)。由于 \(\alpha \in \mathbb{R} - \mathbb{R}_0\) 蕴含 \(\frac{1}{2}\alpha \in \mathbb{R}_0\)，且 \(s_\alpha = s_{\alpha/2}\)，我们有 \(\mathrm{W}(\mathrm{R}) = \mathrm{W}(\mathrm{R}_0)\)。因此 \(\mathbb{R}_0\) 不可约（命题 5 的推论），且显然是约化的。

由于 R 不是约化的，存在 \(\alpha\in R_{0}\)，使得 \(2\alpha\in R\)。由于 \(R_{0}\) 不可约且 \(\dim V\geqslant2\)，\(\alpha\) 不可能与每个根成比例或正交。取 \(\beta\in R_{0}\)，使得 \(n(\beta,\alpha)\neq0\)，且 \(\beta\) 不与 \(\alpha\) 成比例。必要时将 \(\beta\) 换成 \(-\beta\)，可假设 \(n(\beta,\alpha)>0\)。现在 \(\frac{1}{2}n(\beta,\alpha)=n(\beta,2\alpha)\in\mathbf{Z}\)，所以 \(n(\beta,\alpha)\in2\mathbf{Z}\)。由第 3 号的列表，\(n(\beta,\alpha)=2\)，\((\beta|\beta)=2(\alpha|\alpha)\)。由于 \(R_{0}\) 约化，命题 12 表明，对所有 \(\gamma\in R_{0}\)，要么 \((\gamma|\gamma)=(\alpha|\alpha)\)，要么 \((\gamma|\gamma)=2(\alpha|\alpha)\)。此外，上面表明，对所有 \(\gamma\in R-R_{0}\)，向量 \(\frac{1}{2}\gamma\) 是 \(R_{0}\) 的元素，且 \(\left(\frac{1}{2}\gamma|\frac{1}{2}\gamma\right)=(\alpha|\alpha)\)。因此，\(\lambda=(\alpha|\alpha)\)，\(B\neq\varnothing\)，\(R_{0}=A\cup B\)，且 \(R\subseteq A\cup B\cup2A\)；另一方面，若 \(\gamma\in A\)，则存在 \(g\in W(R)\)，使得 \(\gamma=g(\alpha)\)（命题 11），所以 \(2\gamma=g(2\alpha)\in R\)；于是 \(2A\subseteq R\)，且 \(R=A\cup B\cup2A\)。最后，设 \(\gamma,\gamma'\) 是 A 中两个不成比例的元素。于是

\[
n (2 \gamma , \gamma^ {\prime}) = 2 n (\gamma , \gamma^ {\prime}) = 4 n (\gamma , 2 \gamma^ {\prime}) \in 4 \mathbf {Z}, \text { and } | n (\gamma , \gamma^ {\prime}) | \leqslant 1
\]

因为 \(\gamma\) 和 \(\gamma'\) 长度相同，所以 \(n(\gamma, \gamma') = 0\)，且 \((\gamma|\gamma') = 0\)。

命题 14. 假设 R 不可约且约化，并且 \((\alpha|\alpha)\) 取值为 \(\lambda\) 和 \(2\lambda\)，其中 \(\alpha\in R\)。设 A 是根 \(\alpha\) 的集合，满足 \((\alpha|\alpha)=\lambda\)。假设 A 中任意两个不成比例的元素都正交。则 \(R_{1}=R\cup2A\) 是不可约的非约化根系，而 R 是 \(R_{1}\) 的不可分根的集合。

显然，\(R_{1}\) 满足 \((RS_{I})\) 和 \((RS_{III})\)。我们证明，若 \(\alpha, \beta \in R_{1}\)，则 \(\langle\alpha^{\sim}, \beta\rangle \in Z\)。当 \(\alpha \in R\) 时这显然成立。由于 \((2\alpha)^{\sim} = \frac{1}{2}\alpha^{\sim}\)（\(\alpha \in A\)），当 \(\alpha, \beta \in 2A\) 时也立即可得。最后，设 \(\beta \in R\)，且 \(\alpha = 2\gamma\)，其中 \(\gamma \in A\)。

\begin{enumerate}
\def\labelenumi{\arabic{enumi})}
\item
  若 \(\gamma = \pm \beta\)，则 \(\langle \alpha^{\prime},\beta \rangle = \pm \frac{1}{2}\langle \gamma^{\prime},\gamma \rangle = \pm 1\)。
\item
  若 \(\gamma\) 不与 \(\beta\) 成比例且 \(\beta \in A\)，根据 \(A\) 的假设，有 \(\langle \gamma^{\sim},\beta \rangle = 0\)，所以 \(\langle \alpha^{\sim},\beta \rangle = 0\)。
\item
  若 \(\beta \in \mathbb{R} - \mathbb{A}\)，则 \((\beta |\beta) = 2\lambda = 2(\gamma |\gamma)\)，所以根据第 3 号的列表，\(\langle \beta, \gamma^{\prime} \rangle\) 等于 0、2 或 -2。因此 \(\langle \beta, \alpha^{\prime} \rangle = \frac{1}{2}\langle \beta, \gamma^{\prime} \rangle \in \mathbf{Z}\)。
\end{enumerate}

因此 \(R_{1}\) 是 V 中的根系，其余断言显然。

